% !TeX program = pdflatex
\documentclass[11pt,a4paper]{article}
\usepackage[T1]{fontenc}
\usepackage[margin=27mm]{geometry}
\usepackage{amsmath,amssymb,amsthm,mathtools}
\usepackage{booktabs,longtable,array,enumitem}
\usepackage[colorlinks=true,linkcolor=blue,citecolor=blue,urlcolor=blue]{hyperref}
\usepackage{microtype}
\setlength{\parskip}{0.25em}
\setlist{nosep}
\numberwithin{equation}{section}
\newtheorem{theorem}{Theorem}[section]
\newtheorem{proposition}[theorem]{Proposition}
\newtheorem{lemma}[theorem]{Lemma}
\newtheorem{corollary}[theorem]{Corollary}
\theoremstyle{definition}
\newtheorem{definition}[theorem]{Definition}
\theoremstyle{remark}
\newtheorem{remark}[theorem]{Remark}
\newcommand{\R}{\mathbb R}
\newcommand{\dd}{\,\mathrm d}
\newcommand{\dx}{\dd x}
\newcommand{\loc}{\mathrm{loc}}
\newcommand{\SO}{\mathrm{SO}}
\newcommand{\BMO}{\mathrm{BMO}}
\newcommand{\supp}{\operatorname{supp}}
\newcommand{\diver}{\operatorname{div}}
\newcommand{\fint}{\mathop{\int\!\!\!\!-}}
\newcommand{\M}{\mathcal M}
\newcommand{\Th}{\Theta}
\newcommand{\norm}[1]{\left\lVert#1\right\rVert}
\newcommand{\pair}[2]{\left\langle#1,#2\right\rangle}
\newcommand{\cl}[1]{\overline{#1}}
\urlstyle{tt}
\newcommand{\lean}[1]{\nolinkurl{#1}}
\title{Small-energy regularity for stationary harmonic maps\\
\large A complete proof for compact embedded $C^2$ targets}
\author{}
\date{}
\begin{document}
\maketitle
\begin{abstract}
Let $N$ be a nonempty compact embedded $C^2$ submanifold without boundary
in Euclidean space. In every domain dimension $n\geq2$, a stationary weakly
harmonic map whose normalized energy is sufficiently small on a ball has a
unique continuous $N$-valued representative on the concentric half ball.
We prove the scale-correct estimate with H\"older exponent $1/2$ and the
square root of the original normalized energy. The proof uses the normal
projection antisymmetric system, a natural Coulomb frame with critical
Morrey control, compensation for the actual bounded cutoff of the
compact-target map, and an all-center decay iteration. The frame is
constructed by stronger Morrey perturbations, boundary absorption, and a
joint critical limit. The Hodge fields are selected once and retained
through the estimates. No minimizing property and no separately assumed
gauge, decay estimate, or regular representative enter the theorem. A final
correspondence table identifies the machine-checked proof components.
\end{abstract}
\tableofcontents

\section{The exact statement and weak conventions}\label{sec:statement}
Fix integers $n\geq2$, $q\geq0$, and a nonempty compact embedded $C^2$
submanifold $N\subset\R^q$ without boundary. No orientation, global normal
frame, or smoothness beyond $C^2$ is required. All matrix and gradient
norms are Hilbert--Schmidt norms. A weak gradient is the distributional
gradient of the original map, and is kept fixed throughout. Set
\[
 E(u;V)=\int_V|Du|^2\dx,\qquad
 \Th_u(a,r)=r^{2-n}E(u;B_r(a)).
\]
There is no factor $1/2$ in the energy.
For $y\in N$, let $Q(y)$ be the orthogonal projection onto $(T_yN)^\perp$
and $\Pi(y)=I-Q(y)$. The map $Q$ is $C^1$ along $N$ and admits a compactly
supported $C^1$ ambient extension $\widehat Q$ with bounded derivative.
Use the sign convention
$\mathcal A_y(v,w)=DQ(y)[v]w$ for tangent vectors $v,w$.

\begin{definition}\label{def:weak}
A map $u\in W^{1,2}_{\loc}(U,\R^q)$, with $u\in N$ almost everywhere, is
weakly harmonic if the following holds for every
$V\in C_c^1(\R^n\times\R^q,\R^q)$ whose support is contained in
$U\times\R^q$ and for which $V(x,y)\in T_yN$ for every $x\in\R^n$,
$y\in N$: the integrand below is integrable and
\begin{equation}\label{eq:map-variation}
 \int_U\sum_i\pair{\partial_i u(x)}
 {DV(x,u(x))[e_i,\partial_i u(x)]}\dx=0.
\end{equation}
This is the complete ambient map-variation condition. In particular,
after a compact ambient cutoff equal to one near $N$, it includes the
variation $V(x,y)=\Pi(y)\varphi(x)$ for every compactly supported
ambient $C^1$ field $\varphi$. The resulting weak equation is derived
in Lemma~\ref{lem:geometry}. The map is domain stationary if, in addition,
\begin{equation}\label{eq:stress}
 \int_U\left(|Du|^2\diver X
 -2\sum_{i,j=1}^n\pair{\partial_i u}{\partial_j u}\partial_iX_j\right)\dx=0
 \quad(X\in C_c^\infty(U,\R^n)).
\end{equation}
\end{definition}

\begin{theorem}[Small-energy regularity]\label{thm:main}
There are $\varepsilon_*>0$ and $C<\infty$, depending only on $n$ and the
fixed embedded target $N\subset\R^q$, with the following property.
Suppose $U\subset\R^n$ is open, $u\in W^{1,2}_{\loc}(U,\R^q)$ takes
values in $N$ almost everywhere, and $u$ is weakly harmonic and domain
stationary in the sense of Definition~\ref{def:weak}. If $R>0$,
$\cl{B_R(a)}\subset U$, and
\begin{equation}\label{eq:small-energy}
 \Th_u(a,R)\leq\varepsilon_*,
\end{equation}
there is a continuous representative
$\widetilde u:B_{R/2}(a)\longrightarrow N$ satisfying
\begin{equation}\label{eq:holder-main}
 |\widetilde u(x)-\widetilde u(y)|
 \leq C R^{-1/2}\Th_u(a,R)^{1/2}|x-y|^{1/2}
 \qquad(x,y\in B_{R/2}(a)).
\end{equation}
It is unique among all continuous $\R^q$-valued representatives of $u$
on the half ball. The constants precede the domain, map, center, and radius.
\end{theorem}

The argument follows the antisymmetric-potential and Morrey-frame method
of Riv\`iere and Struwe \cite{R,RS}. Its analytic application here is to the
compact-target map itself. We do not assert the wider regularity theorem
for arbitrary, possibly unbounded, Euclidean-valued solutions of
antisymmetric systems. Compactness supplies a genuine bound for a cutoff
used in the native compensation pairings; the estimates do not depend on
the size of this bound.

The accompanying Lean theorem has precisely the stated hypotheses, with
the original explicit weak gradient. Its proof has no unformalized
analytic assumption. The prose uses ordinary background real analysis
(Lebesgue differentiation, weak Sobolev calculus, Hilbert projection, and
the Banach inverse function theorem) and gives the analytic constructions
needed here. Section~\ref{sec:formal} identifies their formal counterparts.

\section{Geometry and stationarity}\label{sec:geometry}
\begin{lemma}[Tangency and the original weak equation]\label{lem:geometry}
One has $Q(u)\partial_i u=0$ almost everywhere and
\begin{equation}\label{eq:harmonic-equation}
 \int_U Du:D\varphi
 =\int_U\sum_i\mathcal A_u(\partial_i u,\partial_i u)\cdot\varphi
 \quad(\varphi\in C_c^\infty(U,\R^q)).
\end{equation}
The skew one-form
\begin{equation}\label{eq:omega}
 \Omega_i=Q(u)\partial_i(Q\circ u)-\partial_i(Q\circ u)Q(u)
\end{equation}
satisfies
\begin{equation}\label{eq:antisym-system}
 -\Delta u=\sum_i\Omega_i\partial_i u,
 \qquad |\Omega|\leq C_N|Du|.
\end{equation}
\end{lemma}
\begin{proof}
Local defining functions for the embedded manifold and the $C^1$ Sobolev
chain rule show that $Du$ is tangent almost everywhere. The same rule gives
$D(\widehat Q\circ u)=D\widehat Q(u)Du$. Use the tangent test
$\Pi(u)\varphi$ in the first variation. Because $\Pi(u)Du=Du$,
$D\Pi=-DQ$, and $DQ$ is bounded on the compact target, the product rule
gives
\[
 0=\int_U Du:D(\Pi(u)\varphi)
  =\int_U Du:D\varphi
     -\int_U\sum_i DQ(u)[\partial_i u]\partial_i u\cdot\varphi.
\]
These bounded native tests are legitimate: bounded smooth approximations
of their weak graphs converge strongly in $W^{1,2}$, with a uniform value
bound obtained by Lipschitz truncation. The derivative terms converge by
Holder, and the $L^1$ forcing terms by dominated convergence. This proves
the equation from the original variations.

Both $Q(u)$ and its derivatives are symmetric; their commutator is skew.
Differentiating $Q^2=Q$ shows that $DQ(y)[v]w$ is normal for tangent
$v,w$. Therefore
$\Omega_i\partial_i u=Q(u)\partial_i(Q\circ u)\partial_i u
=\partial_i(Q\circ u)\partial_i u$.
The bound for $D\widehat Q$ and the preceding weak equation complete
the proof, with the sign displayed in \eqref{eq:antisym-system}.
\end{proof}

\begin{lemma}[Exact monotonicity]\label{lem:monotonicity}
If $0<r<s$ and $\cl{B_s(a)}\subset U$, then
\begin{equation}\label{eq:monotonicity}
 \Th_u(a,s)-\Th_u(a,r)
 =2\int_{B_s(a)\setminus B_r(a)}
 |x-a|^{2-n}|\partial_{\mathrm{rad},a}u|^2\dx\geq0,
\end{equation}
where $\partial_{\mathrm{rad},a}u=Du(x)(x-a)/|x-a|$ away from $a$.
\end{lemma}
\begin{proof}
Translate $a$ to zero and put $X(x)=\eta(|x|)x$ in \eqref{eq:stress},
first with $\eta$ smooth and constant near zero. This gives
\[
 \int\left((n-2)\eta|Du|^2
 +|x|\eta'\bigl(|Du|^2-2|\partial_{\rm rad}u|^2\bigr)\right)=0.
\]
For $F(t)=\int_{B_t}|Du|^2$, polar integration and radial approximation
yield, for almost every $t>0$,
\[
 (n-2)F(t)=tF'(t)-2t\int_{\partial B_t}|\partial_{\rm rad}u|^2
                                      \dd\mathcal H^{n-1}.
\]
$F$ is absolutely continuous on every compact positive radius interval.
Differentiate $t^{2-n}F(t)$, integrate, and use polar coordinates again.
This proves the equality for every positive pair of radii, also for
$n=2$. No minimizing comparison has been used.
\end{proof}

For a field on a ball $D$, use the complete truncated critical quantity
\begin{equation}\label{eq:morrey}
 \M_2(F;D)^2=
 \sup_{x\in D,\,t>0}t^{2-n}\int_{B_t(x)\cap D}|F|^2\dx.
\end{equation}
\begin{lemma}[All-center control from one energy]\label{lem:local-morrey}
If $e=\Th_u(a,R)$ and $D=B_{3R/4}(a)$, then
\begin{equation}\label{eq:joint-smallness}
 \M_2(Du;D)^2+\M_2(\Omega;D)^2\leq M_{n,N}e.
\end{equation}
\end{lemma}
\begin{proof}
For $x\in D$ and $t\leq R/4$, monotonicity at $x$ between $t$ and
$R/4$, and $B_{R/4}(x)\subset B_R(a)$, give
$t^{2-n}\int_{B_t(x)\cap D}|Du|^2\leq4^{n-2}e$.
For $t>R/4$, bound $t^{2-n}$ by $(R/4)^{2-n}$ and use the full
energy on $B_R(a)$. Finally $|\Omega|\leq C_N|Du|$.
\end{proof}

\section{Ball projections and native compensation}\label{sec:analysis}
An interior divergence-free $L^2$ field $Z$ satisfies
$\int Z:D\psi=0$ for compactly supported tests. A \emph{natural}
divergence-free field satisfies this for every $\psi\in W^{1,2}(D)$,
including its weak zero normal flux. The distinction matters below.

\subsection{One Dirichlet and natural decomposition}\label{sec:projection}
\begin{lemma}[Ball projections]\label{lem:projection}
Let $D=B_r(a)$ and $F\in L^2(D,\R^q\otimes\Lambda^1)$. There are
$f\in W^{1,2}_0(D,\R^q)$ and a mean-zero $w\in W^{1,2}(D,\R^q)$
uniquely characterized by
\begin{align}
 \int_D Df:D\varphi&=\int_D F:D\varphi
                &&(\varphi\in W^{1,2}_0(D)),\label{eq:dir-proj}\\
 \int_D Dw:D\psi&=\int_D F:D\psi
                &&(\psi\in W^{1,2}(D)).\label{eq:neu-proj}
\end{align}
Set $Q=F-Dw$, $h=Dw-Df$. Then
\begin{equation}\label{eq:decomposition}
 F=Df+Q+h,
\end{equation}
$Q$ is natural divergence-free, $h$ is closed and divergence-free with
harmonic coefficients, and the three fields are mutually $L^2$ orthogonal.
Each projection is bounded on every $L^s$, $1<s<\infty$, with constants
depending only on $n,s$. If $p'=p/(p-1)$, then
\begin{equation}\label{eq:dirichlet-duality}
 \norm{Df}_{L^p(D)}\leq C_{n,p}
 \sup_{\substack{\varphi\in W^{1,p'}_0(D)\\\norm{D\varphi}_{p'}\leq1}}
       \left|\int_D F:D\varphi\right|.
\end{equation}
For $0<\theta\leq1/2$ one also has
\begin{equation}\label{eq:projection-small-ball}
 \int_{B_{\theta r}(a)}|F|^p
 \leq C_{n,p}\left(\norm{Df}_p^p+\norm Q_p^p
                    +\theta^n\norm F_p^p\right).
\end{equation}
\end{lemma}
\begin{proof}
Poincare makes the gradient norm coercive on the zero-boundary and
mean-zero spaces. Hilbert projection gives \eqref{eq:dir-proj} and
\eqref{eq:neu-proj}. Subtraction for compactly supported tests gives
$\diver h=0$, and $dh=0$ since $h$ is a difference of gradients.
Thus its coefficients are harmonic. The defining identities give all
three orthogonalities.

Here is the analytic mechanism for the all-power bounds. The gradient
response kernels of the ball Dirichlet and mean-zero Neumann problems
satisfy
\[
 |K(x,y)|\leq C|x-y|^{-n},\qquad
 |K(x,y)-K(x',y)|\leq C|x-x'|\,|x-y|^{-n-1}
\]
when $2|x-x'|\leq|x-y|$. These follow from the Green response and
harmonic gradient estimates, including the matching homogeneous boundary
comparison on a spherical cap. One may obtain the latter by inversion
in the sphere with the normal-flux transport for the natural problem.
Split a source at twice the radius of a test ball. The near response is
controlled by the $L^2$ projection norm; the far response is harmonic
and satisfies the displayed difference bound. The disjoint-ball
decomposition of a source above a height then gives weak $(1,1)$:
its good part uses $L^2$, and the mean-zero bad parts use the difference
bound outside the doubled balls. Interpolation gives $1<s\leq2$.
Self-adjointness supplies $s>2$ by duality. This construction extends
the same $L^2$ operators, rather than choosing new solutions for each $s$.

Apply the Dirichlet projection to the power dual of $Df$, using
truncation and then monotone convergence when necessary. Self-adjointness
and the $L^{p'}$ bound give \eqref{eq:dirichlet-duality}.
Finally harmonic mean-value estimates give
$\norm h_{L^p(B_{\theta r})}\leq C\theta^{n/p}\norm h_{L^p(D)}$.
The projection bound controls the last norm by $C\norm F_p$;
\eqref{eq:decomposition} proves \eqref{eq:projection-small-ball}.
\end{proof}

The natural field $Q$ is equivalently the coexact part $*dg$ with its
relative, co-closed potential on a ball. We use its native projection
characterization, retaining the whole boundary test space.

\subsection{The precise div--curl and moment estimates}\label{sec:compensation}
For scalar $b$, put
$[b]_{\BMO(\R^n)}=\sup_B\fint_B|b-b_B|$. Only BMO functions
obtained from actual Sobolev graphs are required here.

\begin{lemma}[Local native div--curl pairing]\label{lem:divcurl}
Let $D=B_r(a)$, let $Z\in L^2(D,\R^k\otimes\Lambda^1)$ be
interior divergence-free, and $T\in W^{1,2}_0(D,\R^k)$.
If a bounded native $W^{1,2}$ scalar $b$ has a whole-space BMO extension
$\widehat b$ agreeing on $D$, then the native pairing agrees with its
ordinary integral and
\begin{equation}\label{eq:clms}
 \left|\int_D b\,Z:DT\right|
 \leq C_n[\widehat b]_{\BMO(\R^n)}\norm Z_2\norm{DT}_2.
\end{equation}
The constant is independent of the bound for $b$.
The same estimate holds when $Z$ is natural divergence-free and $T$
is any native $W^{1,2}$ graph on $D$, with no zero-boundary condition.
\end{lemma}
\begin{proof}
Extend $DT$ by zero. There is a solenoidal $L^2$ extension $E$ of $Z$
to $\R^n$ with norm at most $C_n\norm Z_2$. To construct it,
transport an extension across the sphere and correct its divergence on
an annulus by solving the mean-zero weak divergence problem there;
the compatibility condition is the original divergence identity.
Equivalently the annular correction is defined by duality against
gradients and Hilbert projection. The correction is outside $D$ and
therefore preserves the original field there.

For smooth approximations and a normalized kernel $\rho_t$ supported
in a ball of radius $t$, cancellation gives
\[
 \rho_t*(E\cdot DT)(x)
 =\int \nabla\rho_t(x-y)\cdot E(y)
                    (T(y)-T_{B_{2t}(x)})\dd y.
\]
Choose $2n/(n+1)<s<2$. Sobolev--Poincare and Holder bound the grand
maximal function of $E\cdot DT$ by
\[
 C\bigl(\mathcal M(|E|^s)\bigr)^{1/s}
   \bigl(\mathcal M(|DT|^s)\bigr)^{1/s}.
\]
The maximal operator is bounded on $L^{2/s}$, so the $H^1$ norm is at
most $C\norm E_2\norm{DT}_2$. This is the div--curl cancellation of
\cite{CLMS}, in the precise $L^2$ case needed here.

Its BMO pairing can be seen directly from atoms: the level sets of
the grand maximal function, decomposed into maximal disjoint balls,
give a sum of mean-zero atoms, each bounded by its reciprocal ball
volume, with the sum of coefficients controlled by the $H^1$ norm.
Subtract the ball mean of $\widehat b$ on each atom to bound its pairing
by $C[\widehat b]_{\BMO}$, and sum. Strong $L^2$ approximation of
the weak graphs passes this estimate to the original fields. Bounded
smooth truncations of the actual $b$ show that the limit is its ordinary
integral. Thus its value bound is used for identification and
integrability, and does not enter the estimate.
For the natural version, extend $Z$ by zero: its full natural identity
makes this extension divergence-free on all of $\R^n$. Extend the
centered graph $T-T_D$ across the sphere and cut it off, with gradient
norm at most $C\norm{DT}_2$ by Poincare. The same full-space proof
then applies, while the integral remains exactly the one on $D$.
\end{proof}

\begin{lemma}[Graph BMO and value moments]\label{lem:moments}
The following estimates hold for fixed finite-dimensional Euclidean values.
\begin{enumerate}[label=\textup{(\roman*)}]
\item If $u\in W^{1,2}(B_1)$ and $\M_2(Du;B_1)\leq\eta$, then
\begin{equation}\label{eq:critical-moments}
 \norm{u-u_{B_1}}_{L^k(B_1)}\leq C_{n,k}\eta\quad(0<k<\infty).
\end{equation}
The centered graph has scalar whole-space BMO extensions with seminorm
at most $C_n\eta$ for unit scalar observables.
\item For compactly supported $v\in W^{1,2}(\R^n)$ and $1<p<2$, let
\begin{equation}\label{eq:envelope}
 \Phi(x,r)=r^{p-n}\int_{B_r(x)}|Dv|^p,\qquad
 \Psi(t)=\sup_{x\in\R^n,\,0<r\leq t}\Phi(x,r).
\end{equation}
On $D=B_{t/2}(a)$ every unit scalar observable of $v$, after a common
centering, has a whole-space extension with BMO seminorm at most
$C_{n,p}\Psi(8t)^{1/p}$. If $z=v_{B_{2t}(a)}$, then
\begin{equation}\label{eq:local-moments}
 \norm{v-z}_{L^k(B_t(a))}
 \leq C_{n,p,k}\Psi(8t)^{1/p}t^{n/k}\quad(0<k<\infty).
\end{equation}
\end{enumerate}
\end{lemma}
\begin{proof}
Poincare and Holder bound the mean oscillation on each interior ball
by $C\eta$. To include the boundary, extend the original weak graph
evenly by inversion $x\mapsto x/|x|^2$ from the unit ball to the double
ball. The derivative and Jacobian sizes are $|x|^{-2}$ and $|x|^{-2n}$;
change of variables and recentering intersecting balls show that the
complete truncated critical bound increases by a dimensional factor
only. The two graphs have the same trace on the sphere, hence glue to
a weak graph. Cut off its centered version on the double ball.
Poincare controls the cutoff derivative and gives the global extension.

For clarity, the value moment argument requires only this graph BMO
instance. Decompose a ball into maximal subballs where the local
oscillation exceeds $A[b]_{\BMO}$, for fixed sufficiently large $A$.
Their combined volume is at most half the parent volume. Off these
balls differentiation bounds the oscillation by $C A[b]_{\BMO}$;
their means differ from the parent mean by the same bound. Repeating
gives
$|\{|b-b_B|>jC A[b]_{\BMO}\}\cap B|\leq2^{-j}|B|$.
Integration of $k\lambda^{k-1}$ against this distribution function
proves all finite moments. Centering by $u_{B_1}$ gives (i).

For (ii), rescale $B_{4t}(a)$. Poincare and the actual envelope bound
every subball oscillation by $C\Psi(8t)^{1/p}$. Apply the same graph
extension and cutoff, or use the globally defined graph with a cutoff
on $B_{4t}(a)$. The distribution estimate just proved gives the moments.
Changing a comparable centering mean to $v_{B_{2t}(a)}$ costs only a
fixed factor. These are Sobolev-graph instances; no extension theorem
for an arbitrary abstract ball-BMO function is asserted.
\end{proof}

\section{Construction of the natural Coulomb frame}\label{sec:gauge}
\begin{proposition}[Critical natural frame]\label{prop:gauge}
For fixed $n\geq2,q$, there are $\kappa>0,C_g<\infty$ such that
every skew one-form $A\in L^2(B_1)$ with $\M_2(A;B_1)\leq\kappa$
admits $P\in W^{1,2}(B_1,\SO(q))$ for which
\begin{equation}\label{eq:Z}
 Z_i=P^T\partial_iP+P^TA_iP
\end{equation}
is natural divergence-free and
\begin{equation}\label{eq:gauge-bound}
 \M_2(DP;B_1)+\M_2(Z;B_1)\leq C_g\M_2(A;B_1).
\end{equation}
Thus $Z$ is the literal transformed original connection, not an
independently chosen field. The constants also work on every positive
ball by affine transport.
\end{proposition}

For $0\leq\lambda<n$ write
\[
 \norm F_{2,\lambda}^2=
 \sup_{x\in B_1,\,r>0}r^{-\lambda}
                         \int_{B_r(x)\cap B_1}|F|^2.
\]
The critical exponent is $n-2$; a stronger exponent is
$\nu=n-2+\gamma$, $0<\gamma<2$. We first prove the stronger-class
construction and then take a critical limit.

\begin{lemma}[Boundary estimates and the actual inverse]\label{lem:gauge-operators}
The ball Dirichlet, natural-gradient, and relative coexact responses are
bounded on $L^{2,\lambda}$ for $0\leq\lambda<n$.
If a native orthogonal frame has the natural connection \eqref{eq:Z}
and $\M_2(DP)\leq d_0$, then, for $\lambda=n-2,\nu$ for which
the original quantities $\norm A_{2,\lambda}$ and
$\norm{DP}_{2,\lambda}$ are finite,
\begin{equation}\label{eq:absorption-before}
 \norm{DP}_{2,\lambda}+\norm Z_{2,\lambda}
 \leq C_{n,q,\lambda}\norm A_{2,\lambda}
       +C_{n,q,\lambda}d_0\norm{DP}_{2,\lambda}.
\end{equation}
Let $\mathcal N$ denote the native mean-zero Neumann response, so
$\int D\mathcal N(F):D\psi=\int F:D\psi$ for every $W^{1,2}$ test.
For natural skew $Z$ and a mean-zero native $H$ in the stronger class,
\begin{equation}\label{eq:linearized-neumann}
 \norm{\mathcal N([Z,H])}_{W^{1,2,\nu}}
 \leq C_{n,q,\gamma}\M_2(Z)\norm H_{W^{1,2,\nu}}.
\end{equation}
The norm includes the native gradient and the mean-zero value norm.
\end{lemma}
\begin{proof}
For a linear response $TF$, split the source at twice the radius of a
ball or spherical boundary cap. The near response has energy bounded
by $C\int_{B_{2r}\cap B_1}|F|^2$. The far response is harmonic there
with its matching homogeneous boundary condition. Gradient comparison
on the cap gives the factor $C(r/s)^n$ between its energies at scales
$r$ and $s$. Decomposing into annuli of radii $2^jr$ therefore bounds
the Morrey norm by a geometric series with ratio
$2^{-(n-\lambda)/2}<1$. Radii comparable to or larger than one use
the global energy estimate. This proves the complete truncated-ball
bound, including the boundary. Poincare controls values after the
mean-zero or relative zero-kernel normalization.

Here is a native construction of the nonlinear bound that retains
the full boundary condition throughout. Let $\mathcal Q$ be the
natural solenoidal projection from Lemma~\ref{lem:projection}. Define
\[
 B(X,Y)=\tfrac12(X^TY-Y^TX),\qquad S=P^TdP=B(P,dP).
\]
The last equality is the differentiated orthogonality identity and
$B(X,X)=0$. Let $Q_b=\mathcal Q S$ and
$Q_c=\mathcal Q(P^TAP)$. Since $Z$ is natural, $Z=Q_b+Q_c$.
The linear response bound controls
$\norm{Q_c}_{2,\lambda}\leq C\norm A_{2,\lambda}$.
The notation $Q_b$ refers directly to a one-form; if a potential is
desired, it is $*db$ for the unique relative co-closed $(n-2)$-form
$b$ (a zero-boundary scalar when $n=2$).

For any bounded original critical graph $P$ and any native graph $W$,
let $Q_W=\mathcal Q(B(P,dW))$. Natural orthogonality and the natural
version of Lemma~\ref{lem:divcurl} give
\begin{equation}\label{eq:nonlinear-full-energy}
 \norm{Q_W}_2^2=\int_{B_1}B(P,dW):Q_W
 \leq C[P]_{\BMO}\norm{DW}_2\norm{Q_W}_2
 \leq Cd_0\norm{DW}_2\norm{Q_W}_2.
\end{equation}
The scalar BMO coefficient here is the actual bounded graph $P$.
The natural identity permits zero extension of $Q_W$, so the pairing
uses every whole-ball test and has no residual normal boundary term.
In particular,
\begin{equation}\label{eq:nonlinear-global}
 \int_{B_1}|Q_b|^2\leq Cd_0^2\int_{B_1}|DP|^2.
\end{equation}

To obtain power growth, fix $x\in B_1$, $0<r\leq1/8$, and put
$K=\norm{DP}_{2,\lambda}^2<\infty$. Even inversion extends the
original graph to $B_2$ with both its critical bound and its
$r^\lambda$ energy growth multiplied by fixed constants. Choose
$\chi=1$ on $B_r(x)$, supported in $B_{2r}(x)$, with
$|D\chi|\leq C/r$, and let $c$ be the extended mean on $B_{4r}(x)$.
The restricted native graph $W=\chi(P-c)$ has
$\norm{DW}_{L^2(B_1)}^2\leq CKr^\lambda$ by Poincare on the extended
ball. Alternation eliminates the quadratic cutoff term:
\[
 B(P-c,DW)=\chi B(P-c,DP),\qquad B(P-c,P-c)=0.
\]
Let $Q_{\rm near}=\mathcal Q(\chi B(P-c,DP))$. Applying
\eqref{eq:nonlinear-full-energy} with $P-c$ and $W$ proves
$\int|Q_{\rm near}|^2\leq Cd_0^2Kr^\lambda$.
Subtracting $c$ changes the full forcing only by the actual gradient
$B(c,DP)=D(B(c,P))$, which the same natural projection annihilates.
Consequently $Q_{\rm far}=Q_b-Q_{\rm near}$ is the natural response to
$(1-\chi)B(P-c,DP)$, a forcing that vanishes on $B_r(x)\cap B_1$.

The native gradient complement of this far response is therefore
weakly harmonic there, with homogeneous Neumann boundary data on the
original sphere. Its full natural identity, rather than interior
divergence alone, supplies that boundary condition. The boundary
harmonic gradient estimate, or the native Neumann inversion extension,
gives on $B_{r/256}(x)\cap B_1$
\[
 |Q_{\rm far}|^2\leq Cr^{-n}
      \int_{B_r(x)\cap B_1}|Q_{\rm far}|^2
 \quad\text{almost everywhere}.
\]
Put $E_b(r)=\int_{B_r(x)\cap B_1}|Q_b|^2$. Integrate this estimate
on $B_{\theta r}(x)\cap B_1$, with $0<\theta\leq1/256$, and use
the near-field energy to obtain
\begin{equation}\label{eq:nonlinear-growth-step}
 E_b(\theta r)\leq C\theta^nE_b(r)+Cd_0^2Kr^\lambda.
\end{equation}
Choose $\theta$ so that $C\theta^n\leq\theta^\lambda/2$.
At $r_0=1/8$, \eqref{eq:nonlinear-global} and
$\int_{B_1}|DP|^2\leq K$ give the initial bound
$E_b(r_0)\leq Cd_0^2K$. Iterating
\eqref{eq:nonlinear-growth-step}, and filling radii by monotonicity,
gives $E_b(r)\leq Cd_0^2Kr^\lambda$ at every radius. For large radii
use \eqref{eq:nonlinear-global}; recentering an intersecting ball
controls any ambient center. Thus
$\norm{Q_b}_{2,\lambda}\leq Cd_0\norm{DP}_{2,\lambda}$.
Together with $Z=Q_b+Q_c$ and $DP=PZ-AP$, this proves
\eqref{eq:absorption-before}. At the stronger exponent the assumed
finite original growth is essential; no term with infinite norm is
absorbed to manufacture stronger-class membership.

For \eqref{eq:linearized-neumann}, the critical coefficient is paired
with an oscillation of $H$, rather than its absolute value. Indeed
$\mathcal N([Z,H_D])=0$ for every constant matrix $H_D$, precisely
because $Z$ is natural. On a ball of radius $r$, the stronger graph
has oscillation at most $C\norm H_{W^{1,2,\nu}}r^{\gamma/2}$ by
the dyadic Poincare estimate. Multiplying by the critical energy of
$Z$ gives growth $r^{n-2+\gamma}$. Near--far summation, using the
same constant cancellation on each annulus, controls the mean-zero
response in this stronger Morrey norm. Its value bound follows by
Poincare. This proves \eqref{eq:linearized-neumann}. All graph
identities are obtained by bounded smooth approximation and strong
$L^2$ passage; the response remains the same unique native solution.
\end{proof}

\begin{proof}[Proof of Proposition~\ref{prop:gauge}]
First suppose $A\in L^{2,\nu}$. Let $Y_\nu$ be the full native
Morrey--Sobolev space of matrix-valued maps, with norm the sum of
the size of the mean and $\norm{DH}_{2,\nu}$. It is complete and embeds
into $C^{0,\gamma/2}$: the dyadic mean estimate follows from
$\int_{B_r}|DH|^2\leq Cr^{n-2+\gamma}$. The product rule and the
resulting supremum bound make $Y_\nu$ a unital Banach algebra for
matrix multiplication. Let $X_\nu$ be its closed real linear subspace
of mean-zero skew maps. This subspace is the parameter space, not an
algebra. The matrix exponential is analytic in $Y_\nu$, agrees with
the pointwise exponential, and has its actual differentiated weak graph.

Suppose a stronger-class frame is known for $tA$, with connection $Z$.
For a skew mean-zero $H\in X_\nu$ and a small real $s$, the frame
$P\exp H$ for $(t+s)A$ has connection
\[
 \mathcal F(H,s)=e^{-H}d(e^H)+e^{-H}(Z+sP^TAP)e^H.
\]
The natural equation is $\mathcal N(\mathcal F(H,s))=0$. At $(0,0)$
it vanishes and its derivative in $H$, on the skew mean-zero space, is
$I+\mathcal N([Z,\cdot])$. By
\eqref{eq:linearized-neumann}, the second term has norm less than
$1/2$ if $\M_2(Z)$ is sufficiently small. Its inverse is the
convergent Neumann series. The Banach inverse function theorem gives
actual nearby stronger-class frames. Exponentiating a real skew matrix
preserves special orthogonality and the literal transformed connection.

Choose $d_0$ so the last term of \eqref{eq:absorption-before} can be
absorbed at both exponents. The improved estimates are
\begin{equation}\label{eq:two-exponent}
 \norm{DP}_{2,\lambda}+\norm Z_{2,\lambda}
 \leq C\norm{tA}_{2,\lambda}\qquad(\lambda=n-2,\nu).
\end{equation}
Make the critical size of $A$ sufficiently small that these improved
bounds lie strictly inside the perturbation neighborhood. At $t=0$
use $P=I$. The set of $t\in[0,1]$ with a frame and fixed critical
and stronger bounds from \eqref{eq:two-exponent} is relatively open
by the perturbation argument and the improved bounds. It is relatively
closed by weak $W^{1,2}$ compactness, strong $L^2$ compactness of the
bounded orthogonal frames, and lower semicontinuity on every truncated
ball. Strong convergence of the frames and weak convergence of their
derivatives identify the connection when tested against bounded smooth
fields. The natural identity then passes to all $W^{1,2}$ tests by
density. The stronger growth bounds pass as well. Connectedness of
$[0,1]$ gives the frame at $t=1$.

For critical $A$, extend it by zero and mollify, preserving skewness.
The restrictions $A_j$ satisfy
\[
 A_j\longrightarrow A\text{ strongly in }L^2(B_1),\qquad
 \M_2(A_j;B_1)\leq C_n\M_2(A;B_1).
\]
Jensen and Fubini give the bound. A convolution ball whose center falls
outside the unit ball is recentered at a point of intersection, with
at most twice its radius; the complete truncated definition controls
it. Each mollification belongs to a stronger class, and one may fix
$\gamma=1$. Its stronger norm need not be uniform. Reduce $\kappa$
to allow the dimensional convolution constant and construct frames
$P_j$ with uniform critical bounds. A joint subsequence satisfies
\[
 P_j\to P\text{ strongly in }L^2\text{ and almost everywhere},\qquad
 DP_j\rightharpoonup DP\text{ in }L^2.
\]
The closed group $\SO(q)$ retains the values. Boundedness of $P_j$
and strong convergence of $A_j$ identify $P_j^TA_jP_j$; the strong--weak
limit identifies $P_j^TDP_j$. Thus the resulting field is precisely
\eqref{eq:Z}. Its natural identity passes first for bounded smooth
tests and then by density for all native tests. Lower semicontinuity
on each truncated ball proves \eqref{eq:gauge-bound}.
No critical embedding into continuous functions, or approximation in
the critical Morrey norm, is used. Positive affine transport preserves
the critical seminorm and full natural pairing, giving every ball.
\end{proof}

\section{The compact-target system and its actual cutoff}\label{sec:system}
Work first on $B_1$. Fix
\begin{equation}\label{eq:p}
 p=1+\frac1{2n},\qquad p'=\frac p{p-1}=2n+1>n,\qquad\delta=\frac12.
\end{equation}
Suppose the original compact-target map satisfies
$-\Delta u=A\cdot Du$ and $P$ is the constructed frame. Choose
$\eta\in[0,1]$ so that
\begin{equation}\label{eq:eta}
 \M_2(Du;B_1),\quad\M_2(DP;B_1),\quad\M_2(Z;B_1)\leq\eta.
\end{equation}
Choose a fixed smooth cutoff $\chi$, supported in $B_{3/4}$, equal to
one on $B_{2/3}$, with $0\leq\chi\leq1$ and bounded first two
derivatives. Let
\begin{equation}\label{eq:cutoff}
 c=u_{B_1},\qquad v=\chi(u-c),\qquad
 Dv=\chi Du+(u-c)\otimes D\chi,
\end{equation}
extended by zero to $\R^n$. This is the original global weak graph of
$v$. It is bounded because the target is compact.

\begin{lemma}[The transformed equation and error]\label{lem:cutoff}
On $B_1$ put $F_i=P^T\partial_i v$ and $W=\sum_iZ_iF_i$.
Then $F\in L^2$ and
\begin{equation}\label{eq:cutoff-eq}
 -\diver F=W-e,
\end{equation}
where the locally integrable error is the four-term expression
\begin{equation}\label{eq:error}
 e=\sum_i(\partial_iP^T)(u-c)\partial_i\chi
   +2\sum_iP^T\partial_i u\partial_i\chi
   +P^T(u-c)\Delta\chi
   +\sum_i Z_iP^T(u-c)\partial_i\chi.
\end{equation}
For $0<t\leq1/16$ and $B_t(a)\subset B_1$,
\begin{equation}\label{eq:error-L1}
 \norm e_{L^1(B_t(a))}\leq C\eta t^{n-1-\delta}.
\end{equation}
The actual envelope \eqref{eq:envelope} of $v$ is finite and
\begin{equation}\label{eq:initial-envelope}
 0\leq\Psi(t)\leq C\eta^p\qquad(0<t\leq1/16).
\end{equation}
\end{lemma}
\begin{proof}
Orthogonality gives $\partial_iP^T=-P^T(\partial_iP)P^T$. Test the
original system with the bounded native test $P\varphi$. Product
differentiation yields $-\diver(P^TDu)=Z\cdot P^TDu$. Apply the
product rule to \eqref{eq:cutoff} to obtain exactly
\eqref{eq:cutoff-eq}--\eqref{eq:error}. Bounded smooth graph
approximation justifies this for the original weak fields; all terms
on the right are $L^1$.

Choose $k>\max(2,n/\delta)$. The value moments and critical bounds give
\[
 \norm{u-c}_{L^2(B_t(a))}\leq C_k\eta t^{n/2-\delta},\qquad
 \norm{Du}_{L^2(B_t(a))}+\norm{DP}_{L^2(B_t(a))}
 +\norm Z_{L^2(B_t(a))}\leq C\eta t^{n/2-1}.
\]
The first and fourth error terms are at most
$C\eta^2t^{n-1-\delta}$, the second at most $C\eta t^{n-1}$, and
the third at most $C\eta t^{n-\delta}$. Use $\eta,t\leq1$.

Holder gives $\int_{B_t}|Du|^p\leq C\eta^pt^{n-p}$. For the
cutoff derivative take $k>\max(p,n)$; then
$t^{p-n}\int_{B_t}|u-c|^p\leq C\eta^pt^{p-np/k}\leq C\eta^p$.
Any ball of radius at most $1/16$ meeting the support lies in $B_{7/8}$;
otherwise its weak gradient is zero almost everywhere. This proves the
claim for every center in $\R^n$, and finiteness of the supremum.
\end{proof}

\subsection{Compensation for one and the same decomposition}\label{sec:same-fields}
Fix $0<t\leq1/128$ and $D=B_{t/2}(a)\subset B_1$. Apply
Lemma~\ref{lem:projection} once to $F=P^TDv$ on $D$, and retain its
$f,w,Q,h$. Write $b=\Psi(8t)^{1/p}$.

\begin{lemma}[The same two component estimates]\label{lem:field-bounds}
These fields satisfy
\begin{align}
 \norm{Df}_{L^p(D)}&\leq C\eta t^{n/p-1}b+C\eta t^{n/p-\delta},
                                                   \label{eq:f-bound}\\
 \norm Q_{L^p(D)}&\leq C\eta t^{n/p-1}b.\label{eq:Q-bound}
\end{align}
\end{lemma}
\begin{proof}
Take $\varphi\in W^{1,p'}_0(D)$ with $\norm{D\varphi}_{p'}\leq1$.
Since $p'>n$,
\begin{equation}\label{eq:test-scaling}
 \norm\varphi_\infty\leq Ct^{1-n/p'},\qquad
 \norm{D\varphi}_2\leq Ct^{n/2-n/p'}.
\end{equation}
The same Dirichlet equation gives
$\int_D Df:D\varphi=\int_D(W-e)\cdot\varphi$.
Move the derivative of $v$ in $W$ onto the matrix test constructed
from $P\varphi$. In coordinates its derivative has the two terms
$DP\,\varphi+P\,D\varphi$. The resulting pairings are scalar
observables of $v$ against $Z$ and this derivative. Centering constants
vanish by the divergence identity and zero trace. The product test is
native and zero-boundary since both $P$ and $\varphi$ are bounded.
Lemmas~\ref{lem:divcurl} and \ref{lem:moments} give
\[
 \left|\int_D W\cdot\varphi\right|
 \leq Cb\norm Z_2
       (\norm{DP}_2\norm\varphi_\infty+\norm{D\varphi}_2)
 \leq C\eta t^{n/p-1}b.
\]
The actual bounded cutoff identifies this pairing with the ordinary
integral. The error is bounded by
$C\eta t^{n-1-\delta}t^{1-n/p'}=C\eta t^{n/p-\delta}$.
Dirichlet duality proves \eqref{eq:f-bound}.

For $Q$, let $H\in L^{p'}(D)$ be a one-form of norm at most one and
$T$ its natural solenoidal projection. Then
$\norm T_{p'}\leq C$. Self-adjointness and the full natural identity
give, with $z=v_{B_{2t}(a)}$,
\begin{equation}\label{eq:Q-dual}
 \int_D Q:H=\int_D P^TDv:T
          =-\int_D (v-z)\cdot\sum_i\partial_iP\,T_i.
\end{equation}
For smooth bounded approximations this is integration by parts; the
native product test is justified by boundedness of $P$ and $v$.
The constant term vanishes because the natural identity can test the
frame itself. For general duals use approximation and boundedness of
the same natural projection.
Let $s=2p/(2-p)$; then $1/2+1/p'+1/s=1$. Holder and
\eqref{eq:local-moments} bound the last integral by
\[
 C\norm{DP}_2\norm T_{p'}\norm{v-z}_s
 \leq C\eta t^{n/2-1}bt^{n/s}
 =C\eta t^{n/p-1}b.
\]
All limiting products are integrable by this same three-factor bound.
Duality proves \eqref{eq:Q-bound} for the original selected $Q$;
no field is reselected for a dual test.
\end{proof}

\begin{lemma}[One-ball decay]\label{lem:one-ball}
For $0<\theta\leq1/2$,
\begin{equation}\label{eq:one-ball}
 \int_{B_{\theta t/2}(a)}|Dv|^p
 \leq C\left((\eta^p+\theta^n)\Psi(8t)t^{n-p}
                    +\eta^pt^{n-p\delta}\right).
\end{equation}
\end{lemma}
\begin{proof}
Orthogonality gives $|F|=|Dv|$ almost everywhere. Insert the same
component bounds into \eqref{eq:projection-small-ball}. Use
$(a+b)^p\leq2^p(a^p+b^p)$ and
$\norm F_{L^p(D)}^p\leq\Psi(8t)t^{n-p}$, with $n-p\geq0$.
This proves the displayed inequality directly from the original data.
\end{proof}

\section{All-center iteration and the representative}\label{sec:iteration}
\begin{lemma}[Envelope improvement]\label{lem:envelope}
There are $\eta_*>0,C$ such that $\eta\leq\eta_*$ implies
\begin{equation}\label{eq:power-growth}
 \Psi(t)\leq C\eta^pt^\beta\quad(0<t\leq1/16),
 \qquad\beta=p(1-\delta)=p/2.
\end{equation}
For every $a\in\R^n$ and $r>0$,
\begin{equation}\label{eq:gradient-growth}
 \int_{B_r(a)}|Dv|^p\leq C\eta^pr^{n-p+\beta}
                         =C\eta^pr^{n-p/2}.
\end{equation}
\end{lemma}
\begin{proof}
To pass from one ball to all centers, fix $r\leq\theta t/2$ and set
$t'=2r/\theta\leq t$. If $B_r(a)$ misses the support of $Dv$, its
integral is zero. Otherwise $a$ is within distance $r$ of $B_{3/4}$;
since $t\leq1/128$, the closed ball of radius $t'/2$ at $a$ is contained
in $B_1$. Apply \eqref{eq:one-ball} at $t'$, multiply by
$r^{p-n}=(\theta t'/2)^{p-n}$, and use monotonicity of $\Psi$ and
$(t')^\beta\leq t^\beta$. Taking the supremum over all such $a,r$
gives
\begin{equation}\label{eq:raw-recurrence}
 \Psi(\theta t/2)\leq
 D(\eta^p\theta^{p-n}+\theta^p)\Psi(8t)
     +D\eta^p\theta^{p-n}t^\beta
 \quad(0<t\leq1/128).
\end{equation}
Every supremum is finite by \eqref{eq:initial-envelope}; neither
attainment nor continuity is needed. Set $\sigma=\theta/16$ and
replace $t$ by $s/8$. First choose $\theta\in(0,1/2)$ so small that
$D\theta^p\leq\sigma^\beta/4$, possible since $\beta<p$.
Then choose $\eta_*$ so
$D\eta_*^p\theta^{p-n}\leq\sigma^\beta/4$.
This yields
\begin{equation}\label{eq:recurrence}
 \Psi(\sigma s)\leq\rho\Psi(s)+A\eta^ps^\beta,
 \qquad0\leq\rho\leq\tfrac12\sigma^\beta<\sigma^\beta.
\end{equation}
For $s_j=\sigma^j/16$, iteration gives
\[
 \Psi(s_j)\leq\rho^j\Psi(1/16)
 +A\eta^p(1/16)^\beta
      \sum_{k=0}^{j-1}\rho^{j-1-k}\sigma^{k\beta}
 \leq C\eta^ps_j^\beta.
\]
Monotonicity of $\Psi$ fills the gaps between successive $s_j$,
costing the fixed factor $\sigma^{-\beta}$. Multiplying by $r^{n-p}$
gives \eqref{eq:gradient-growth} for $r\leq1/16$.
For larger $r$, the global cutoff estimate
$r^{p-n}\int_{B_r(a)}|Dv|^p\leq C\eta^p$ follows by the same Holder
bounds as before and compact support for $r\geq1$. Since
$r^\beta\geq(1/16)^\beta$, enlarging $C$ completes the proof.
\end{proof}

\begin{lemma}[Direct mean-limit representative]\label{lem:representative}
Suppose $v\in W^{1,2}(\R^n,\R^q)$, $1\leq p\leq2$, $0<\alpha<1$,
and
$\int_{B_r(a)}|Dv|^p\leq K^pr^{n-p+p\alpha}$ for all $a,r>0$.
Then $v$ has a continuous representative with
$[\widetilde v]_{C^{0,\alpha}}\leq C_{n,p,\alpha}K$.
\end{lemma}
\begin{proof}
Poincare and Holder give
$\fint_{B_r(a)}|v-v_{B_r(a)}|\leq CKr^\alpha$.
Successive means at radii $r,r/2,r/4,\ldots$ differ by at most
$CK(r/2^j)^\alpha$. Their limits define $\widetilde v(a)$ at every
point, with tail at most $CKr^\alpha$. Lebesgue differentiation
identifies the limit with $v$ almost everywhere. For distinct $a,b$,
take $r=2|a-b|$ and compare their radius-$r$ means through a ball
of comparable radius containing both. Poincare bounds the mean
difference and the two tails by $CKr^\alpha$. This proves the
pointwise H\"older estimate and continuity everywhere.
\end{proof}

\begin{proposition}[The compact-target representative]\label{prop:system-rep}
The original system and its constructed frame satisfying
\eqref{eq:eta} with $\eta\leq\eta_*$ have a unique continuous
representative on $B_{2/3}$, with H\"older exponent $1/2$ and seminorm
at most $C\eta$. On a ball $B_\rho(a)$ the corresponding conclusion
is on $B_{2\rho/3}(a)$ with bound $C\eta\rho^{-1/2}$.
\end{proposition}
\begin{proof}
Apply \eqref{eq:gradient-growth} and the preceding mean-limit lemma.
On $B_{2/3}$ the cutoff is one, so $c+\widetilde v$ represents the
original $u$. It lies in $N$ at every point: continuity and closedness
of $N$ would otherwise give an open ball outside $N$, contradicting
almost-everywhere target membership. Two continuous representatives
agree everywhere because a nonzero difference persists on an open set
of positive measure. Under $x=a+\rho y$, the weak gradient has the
factor $\rho$, the critical quantities are invariant, and the H\"older
seminorm has the factor $\rho^{-1/2}$.
\end{proof}

\section{Proof of the small-energy theorem}\label{sec:assembly}
\begin{proof}[Proof of Theorem~\ref{thm:main}]
Put $e=\Th_u(a,R)$ and $\rho=3R/4$. Geometry and stationarity give
the original system on this same ball and
\[
 \M_2(Du;B_\rho(a))^2+\M_2(\Omega;B_\rho(a))^2\leq M_{n,N}e.
\]
Let $\zeta=\sqrt{M_{n,N}}\sqrt e$. Choose $\varepsilon_*>0$ so
that $\zeta\leq\kappa$ and
$(1+C_g)\zeta\leq\min(1,\eta_*)$ whenever $e\leq\varepsilon_*$.
These constants depend only on $n$ and the fixed embedded target.
Construct the natural frame for the original connection $\Omega$ on
$B_\rho(a)$ by Proposition~\ref{prop:gauge}. The three bounds in
\eqref{eq:eta} then hold with $\eta=(1+C_g)\zeta$.

Proposition~\ref{prop:system-rep} applies to the actual original map
and this frame. Its domain of conclusion is exactly
$B_{2\rho/3}(a)=B_{R/2}(a)$ and its estimate is
\[
 [\widetilde u]_{C^{0,1/2}(B_{R/2}(a))}
 \leq C(1+C_g)\sqrt{M_{n,N}}\sqrt e\,(3R/4)^{-1/2}.
\]
Absorb fixed factors into $C$. Target membership and uniqueness come
from the same proposition. The case $e=0$ is included: Poincare gives
the constant representative when the original gradient is zero, and
the right side of the estimate is zero. All constants preceded
$U,u,Du,a,R$. No frame, decay, or representation premise has been added
to the theorem.
\end{proof}

\section{The smooth-target corollary}\label{sec:smooth}
\begin{corollary}\label{cor:smooth}
If $N$ is a compact smooth embedded target, the representative of
Theorem~\ref{thm:main} is smooth on the same half ball. It retains
\eqref{eq:holder-main}, all-point target membership, and uniqueness
among continuous representatives. No quantitative bound for all higher
derivatives is asserted.
\end{corollary}
\begin{proof}
We use the standard interior smooth bootstrap for locally H\"older
weak solutions of a smooth quadratic-gradient Poisson equation:
if $w\in W^{1,2}_{\loc}(V)$ is locally H\"older, $A$ is a smooth
bilinear coefficient depending on $w$, and
\[
 \Delta w=\sum_i A(w)(\partial_iw,\partial_iw)
\]
in the actual weak sense, then $w$ is smooth locally. In this criterion
the initial weak gradient and weak equation are retained. Analytically,
one freezes the coefficient on a sufficiently small ball, uses the
small oscillation of $w$ and harmonic comparison to obtain the first
derivative regularity, and then applies interior elliptic estimates and
Schauder bootstrap to the smooth quadratic coefficient. The statement
is a regularity theorem about the weak equation, and is not the false
claim that an arbitrary continuous quadratic system can be differentiated
classically without this first regularity step. We use this proved
bootstrap rather than reproving general interior elliptic theory here;
its complete formal implementation is the theorem
\lean{HolderToCK.holderQuadraticPoisson_smooth} in
\lean{HolderToCK/Frozen/HolderQuadraticPoissonSmooth.lean}.

For the present map, choose a smooth ambient normal-projection
extension $G$. Lemma~\ref{lem:geometry} gives
$\Delta u=-\sum_i DG(u)[\partial_i u]\partial_i u$.
Almost-everywhere replacement by the H\"older representative preserves
the original weak gradient and the equation on the half ball; the
coefficient becomes $A(\widetilde u)=-DG(\widetilde u)$, which is
smooth in its argument. All hypotheses of the stated bootstrap are
therefore supplied internally by Theorem~\ref{thm:main}, the original
variation condition, and smooth target geometry. The conclusion is
smoothness of this same representative on every interior ball of the
half ball. Continuous uniqueness identifies the local smooth versions.
The original quantitative H\"older estimate is unchanged.
\end{proof}

\section{Correspondence with the formal proof}\label{sec:formal}
The public theorem is
\lean{StationaryHarmonicMaps.stationary_small_energy_regularity}.
It chooses $\varepsilon,\alpha,C$ after the fixed dimension and target,
and before the open domain, original map and weak gradient, center,
and positive radius. The proof chooses $\alpha=1/2$. Its conclusion
covers the entire half ball, all-point target membership, the original
energy square-root factor, and uniqueness among all continuous
almost-everywhere representatives.

The weak-harmonicity premise is the original tangent-variation
condition. The proof internally constructs the normal-projection
system, complete Morrey bounds, and native frame, then supplies every
premise of the compensated representative result. No Hodge component,
compensation bound, recurrence, or regular representative is assumed.
The following paths are relative to the repository root, under
\lean{StationaryHarmonicMaps/}. Intermediate analytic lemmas may have
inputs produced by preceding theorems; all such inputs are supplied
inside the main proof.

\small
\begin{longtable}{@{}p{0.27\textwidth}p{0.67\textwidth}@{}}
\toprule
Mathematical step & Formal source\\\midrule\endhead
Embedded target and original variations &
\lean{Foundation/C2Target.lean};\newline
\lean{Foundation/WeakHarmonicity.lean}\\
Projection geometry and original equation &
\lean{Frozen/NormalProjectionCharacterization.lean};\newline
\lean{Frozen/VariationalWeakPoissonEquivalence.lean}\\
Monotonicity and complete Morrey bound &
\lean{Frozen/Monotonicity.lean};\newline
\lean{Frozen/StationaryMorreyBound.lean};\newline
\lean{Analysis/StationarySystemSmallness.lean}\\
Ball projections and all powers &
\lean{Analysis/NeumannGradientProjection.lean};\newline
\lean{Analysis/NeumannOperatorAllLp.lean};\newline
\lean{Analysis/HodgeLowPComparison.lean}\\
Actual graph BMO and moments &
\lean{Analysis/BoundaryCriticalBMOExtension.lean};\newline
\lean{Analysis/CriticalLpMorreyRealMoments.lean};\newline
\lean{Analysis/CappedEnvelopeValueMoments.lean}\\
Solenoidal extension and native pairing &
\lean{Analysis/SolenoidalBallExtension.lean};\newline
\lean{Analysis/LocalDivCurlHardy.lean};\newline
\lean{Analysis/LocalDivCurlContinuousPairing.lean}\\
Boundary Morrey estimates and absorption &
\lean{Analysis/CoulombBoundaryAbsorption.lean};\newline
\lean{Analysis/CoulombTwoExponentAbsorption.lean}\\
Actual stronger-class perturbation &
\lean{Analysis/NativeCoulombPerturbations.lean};\newline
\lean{Analysis/StrongerNativeCoulombExistence.lean}\\
Joint critical limit and native frame &
\lean{Analysis/CriticalNativeCoulombExistence.lean};\newline
\lean{Analysis/CriticalNativeCoulombBallExistence.lean}\\
Same cutoff and compensated fields &
\lean{Analysis/SystemCutoffCompensation.lean};\newline
\lean{Analysis/SystemCoexactCompensation.lean};\newline
\lean{Analysis/SystemCompensatedHodgeFields.lean}\\
One-ball decay and capped envelope &
\lean{Analysis/OneBallCompensatedDecay.lean};\newline
\lean{Analysis/CappedEnvelopeImprovement.lean};\newline
\lean{Analysis/CompensatedCutoffPowerGrowth.lean}\\
Mean-limit representative and scaling &
\lean{Analysis/LocalLpMorreyRepresentative.lean};\newline
\lean{Analysis/CompensatedSystemRepresentative.lean}\\
Original theorem assembly &
\lean{Providers/SmallEnergyRegularity.lean};\newline
\lean{Frozen/SmallEnergyRegularity.lean}\\
Smooth-target corollary &
\lean{Analysis/SmoothHolderRepresentative.lean};\newline
\lean{Providers/SmallEnergySmoothRegularity.lean};\newline
\lean{Frozen/SmallEnergySmoothRegularity.lean}\\
\bottomrule
\end{longtable}
\normalsize

Lean's kernel checks the proof. Its axiom closure contains only
\lean{propext}, \lean{Classical.choice}, and \lean{Quot.sound}.
There is no custom analytic axiom, \lean{sorryAx}, or native trust
bypass. Imported library theorems and the pinned formal dependency
are themselves formal proofs, rather than unformalized analytic
assumptions. The repository README records the reproducible build
and exact theorem verification.

The scope is this closed small-energy theorem and its actual formal
dependencies. It does not certify a separate arbitrary ball-BMO
extension theorem, an arbitrary unbounded-system theorem, or a
$W^{1,p}_{\loc}$ representative theorem on all possible $p\geq1$.
The actual Sobolev-graph and $W^{1,2}$ instances proved here are the
ones consumed by the main theorem. Citations identify the method,
and do not replace missing formal premises.

\begin{thebibliography}{9}
\bibitem{R} T.~Riv\`iere,
\emph{Conservation laws for conformally invariant variational problems},
Invent. Math. \textbf{168} (2007), 1--22.
\href{https://doi.org/10.1007/s00222-006-0023-0}{doi:10.1007/s00222-006-0023-0}.
\bibitem{RS} T.~Riv\`iere and M.~Struwe,
\emph{Partial regularity for harmonic maps, and related problems},
Comm. Pure Appl. Math. \textbf{61} (2008), 451--463.
\href{https://doi.org/10.1002/cpa.20205}{doi:10.1002/cpa.20205}.
\bibitem{CLMS} R.~Coifman, P.-L.~Lions, Y.~Meyer, and S.~Semmes,
\emph{Compensated compactness and Hardy spaces},
J. Math. Pures Appl. (9) \textbf{72} (1993), 247--286.
\bibitem{U} K.~Uhlenbeck,
\emph{Connections with $L^p$ bounds on curvature},
Comm. Math. Phys. \textbf{83} (1982), 31--42.
\bibitem{MR} Y.~Meyer and T.~Riv\`iere,
\emph{A partial regularity result for a class of stationary Yang--Mills
fields in high dimension}, Rev. Mat. Iberoamericana \textbf{19} (2003),
195--219.
\bibitem{IM} T.~Iwaniec and G.~Martin,
\emph{Geometric Function Theory and Non-linear Analysis},
Oxford Mathematical Monographs, Oxford University Press, 2001.
\end{thebibliography}
\end{document}
